\cotefile{156-6}
\scanpage{10}\mspage{9}
\makeatletter\tagsleft@true\makeatother
de \uline{compatibilité} dans $M$ est une \deleted{\unsure{notion}} relation importante, qui se \unsure{surrajoute|surajoute} aux relations $\underset{M}{\leq}$, $\underset{M}{\ll}$. \unsure{Moralement}, pour récupérer \unsure{l'atelier|le treillis} (à iso unique près) il suffit de \unsure{connaitre|connaître} sur $M$ ces deux relations d'ordre, et la relation de compatibilité.
% slanted
\margin{au moins dans le cas \unsure{modéré}}

Les relations \deleted{\unsure{suivantes|restantes}} \inserted{(compatibilité, disjonction)} sont si importantes, qu'elles méritent chacune une notation lapidaire. Je propose
\[
\deleted{F\parallel G}\qquad F\parallel G \tag*{(1.14)}
\]
pour la disjonction, et
\[
\deleted{F\unsure{\mathrel{\tikz[baseline=0pt,x=1em,y=1em]{\draw (-.38,.62)--(.38,.25)--(-.43,-.10); \draw (.34,.62)--(-.32,.25)--(.38,-.10);}}|\mathrel{\tikz[baseline=0pt,x=1em,y=1em]{\draw (-.44,.26)--(0,.58)--(.44,.26)--(0,-.06)--cycle; \draw (0,.58)--(0,-.06);}}}G}\qquad F\mathrel{\tikz[baseline=0pt,x=1em,y=1em]{\draw (-.44,.26)--(0,.58)--(.44,.26)--(0,-.06)--cycle; \draw (0,.58)--(0,-.06);}} G\quad\text{ou}\quad
\tikz[baseline=(fg.base)]{\node[inner sep=1pt] (fg) {$F\mathrel{\tikz[baseline=0pt,x=1em,y=1em]{\draw (-.38,.62)--(.38,.25)--(-.43,-.10); \draw (.34,.62)--(-.32,.25)--(.38,-.10);}} G$}; \draw (fg.center) ellipse [x radius=2.12em,y radius=1.04em];}
\tag*{(1.15)}
\]
pour la compatibilité (NB $\mathrel{\tikz[baseline=0pt,x=1em,y=1em]{\draw (-.38,.62)--(.38,.25)--(-.43,-.10); \draw (.34,.62)--(-.32,.25)--(.38,-.10);}}$ est plus rapide \unsure{à|:} écrire que $\mathrel{\tikz[baseline=0pt,x=1em,y=1em]{\draw (-.44,.26)--(0,.58)--(.44,.26)--(0,-.06)--cycle; \draw (0,.58)--(0,-.06);}}$, \deleted{$\unsure{\diamond|\mathrel{\tikz[baseline=0pt,x=1em,y=1em]{\draw (-.44,.26)--(0,.58)--(.44,.26)--(0,-.06)--cycle; \draw (0,.58)--(0,-.06);}}}$} mais est moins suggestif de l'idée combinatoire exprimée.

On n'a rien dit pour l'instant au sujet de la relation $\deleted{\unsure{\mathrel{\tikz[baseline=0pt,x=1em,y=1em]{\draw (-.38,.62)--(.38,.25)--(-.43,-.10); \draw (.34,.62)--(-.32,.25)--(.38,-.10);}}|\prec|\preccurlyeq}}\preccurlyeq$ de raffinement sur $\unsure{\mathcal{F}|\mathscr{F}}$. Elle est \unsure{formellement|finalement} de nature plus délicate, et est la notion la plus spécifique dans l'``Algèbre des figures''. La signification géométrique s'éclaire par les \unsure{exemples|exemple}, et les \unsure{axiomes|raisons} que \unsure{ceux-|ceux}%
\scanpage{11}\mspage{10}%
ci inspirent.

\uline{Axiomes}\quad On peut prendre pour \unsure{ens} de base \unsure{ds} une axiomatique des ateliers, soit $\unsure{\mathcal{F}|\mathscr{F}}$, soit $M$, soit $\mathcal{L}$. On trouve des données les plus simples sur $\mathcal{F}$, les \deleted{\unsure{plus}} \unsure{sophistiquées|sophistiqués} sur $\mathcal{L}$. Par contre, les axiomes essentiels sont les moins transparents quand on \unsure{reste|entre} \unsure{plus ou moins} dans $\mathcal{F}$, \unsure{le|les} plus \unsure{transparent|transparents} quand on travaille dans $\mathcal{L}$ (ou peut-être pareil pour $M$ ?).

\uline{Sur $\mathcal{F}$}, on peut prendre comme structure de base, au choix
\[
\left\{\begin{aligned}
 &\text{a)}\ \leq,\ \ll\\
 &\text{b)}\ \leq,\ \preccurlyeq \quad \longleftarrow\\
 &\text{c)}\ \ll,\ \deleted{\unsure{F|\mathcal F}}\quad
 \underbrace{\text{et}\quad\mathcal{F}\longrightarrow\mathrm{P}(\mathcal{F})}_{
 \begin{smallmatrix}
 \text{i.e. la \uline{relation}}\\[-1pt]
 \text{\uline{d'incidence}}\\
 \text{``$F$ est une strate de $G$''}\\
 \text{(i.e. $F\in M$ \uline{et} $F\leq G$)}
 \end{smallmatrix}}
\end{aligned}\right.
\quad
\begin{aligned}
 &\text{(notée }F\longmapsto\widetilde F,\ \text{le}\\
 &\text{\unsure{déployement|déploiement} de }F\\
 &\text{vu comme partie de }\mathcal{F}\text{)}\\
 &\text{\deleted{i.e. la \unsure{description|désignation} \inserted{relation \unsure{de type|d'inc.}}}}
\end{aligned}
\tag*{(1.16)}
\]
\margin{(20. juin) Non, \unsure{par mon nouveau point de vue} ça n'implique plus $\ll$}

\deleted{\unsure{Je}} Avant de préciser ces points, je vais grouper dans un diagramme les six relations fondamentales dans $\mathcal{F}$, \unsure{on écrira|en écrivant}
\[
F\triangleleft G\quad
\underset{\unsure{\mathrm{par}}}{\overset{\unsure{\mathrm{dfn}|\mathrm{déf}}}{\Longleftrightarrow}}\quad
\text{``$F$ incident à $G$''}\quad\text{i.e.}\quad F\leq G\ \text{et}\ F\in M
\tag*{(1.17)}
\]
\scanpage{12}\mspage{11}
\diagram{%
\[
\vcenter{\hbox{\begin{tikzcd}[ampersand replacement=\&,row sep=1.25em,column sep=.2em,
execute at end picture={\draw ([xshift=-.65em,yshift=1em]leq.north west) -- ([xshift=.65em,yshift=.1em]leq.north east) -- ([xshift=.75em,yshift=-.95em]ltlt.south east) -- ([xshift=-.9em,yshift=-.1em]ltlt.south west) -- cycle;}]
\& \& G\triangleleft F \arrow[d,Rightarrow] \& \& \\
G\preccurlyeq F \arrow[dr,Rightarrow] \& \& |[alias=leq]| G\leq F \arrow[dl,Rightarrow] \arrow[dr,Rightarrow] \& \& G\parallel F \arrow[dl,Rightarrow] \\
\& |[alias=ltlt]| G\ll F \& \& G\mathrel{\tikz[baseline=0pt,x=1em,y=1em]{\draw (-.38,.62)--(.38,.25)--(-.43,-.10); \draw (.34,.62)--(-.32,.25)--(.38,-.10);}} F \&
\end{tikzcd}}} \tag*{(1.18)}
\]}

Parmi les six relations, il y a quatre relations d'ordre $\triangleleft$, $\leq$, $\preccurlyeq$, $\ll$, et les deux autres relations $\parallel$ et $\mathrel{\tikz[baseline=0pt,x=1em,y=1em]{\draw (-.38,.62)--(.38,.25)--(-.43,-.10); \draw (.34,.62)--(-.32,.25)--(.38,-.10);}}$ sont symétriques.
% slanted
\margin{Non, $\triangleleft$ n'est pas rel. d'ordre, elle est transitive mais non réflexive, car $F\triangleleft F\Longleftrightarrow F\unsure{\in|\overset{\prime}{\in}}M$ Pour avoir une relation d'ordre il faudrait écrire \deleted{$\unsure{F\triangleleft|F}$} $F\trianglelefteq F\overset{\mathrm{déf}}{\Longleftrightarrow}F\triangleleft F$ \uuline{ou} $F=F$.}
La relation $\mathrel{\tikz[baseline=0pt,x=1em,y=1em]{\draw (-.38,.62)--(.38,.25)--(-.43,-.10); \draw (.34,.62)--(-.32,.25)--(.38,-.10);}}$ est \inserted{de plus} réflexive, mais $\parallel$ ne l'est que si $\unsure{\mathcal{F}|\mathscr{F}}$ est réduit à son seul \unsure{élémt|élément} $\varnothing_{\mathcal{F}}$, car
\[
 F\parallel F \quad\Longleftrightarrow\quad F=\varnothing_{\mathcal{F}}\qquad \text{figure vide.} \tag*{(1.19)}
\]

\uline{Remarques}\quad a) Les relations $\triangleleft$ et $\leq$ se déterminent mutuellement. Connaissant $\leq$, on définit $M\subset\mathcal{F}$ comme
\[
M=\left\{\begin{aligned}
&X\in\mathcal{F}\mid X\
\underset{\text{pour }\leq}{\text{$\mathrm{str}^{\mathrm{t}}$. irréductible}},\
\text{i.e.}\ X\leq\mathop{\mathrm{Sup}}\limits_i F_i\Rightarrow\\
&\hspace{12em}\exists i\ \text{avec }X\leq F_i
\end{aligned}\right\}
\tag*{(1.20)}
\]
puis $\triangleleft$ par (1.17).
\deleted{\unsure{Et|En}} D'autre part, connaissant $\triangleleft$ on définit $M$ par
\[
M=\{X\in\mathcal{F}\mid X\triangleleft X\} \tag*{(1.21)}
\]
puis, pour $F\in\mathcal{F}$, son \unsure{déployement|déploiement}
\[
\widetilde F=\{X\in M\mid X\triangleleft\unsure{\mathcal{F}|F}\} \tag*{(1.22)}
\]
et enfin
\[
F\leq G\quad\Longleftrightarrow\quad\widetilde F\subset\widetilde G \tag*{\unsure{(1,23)|(1.23)}}
\]
Donc, pour déterminer \deleted{la structure, il est \unsure{équivalent}} \inserted{les données \unsure{minima|minimes} qui} fixent la structure d'atelier, il est indifférent d'y
\scanpage{13}\mspage{12}
\deleted{\unsure{pr|pre}}inclure soit $\leq$, soit $\triangleleft$.

\deleted{b) $\leq$ et $\ll$ \unsure{déterminent} $\preccurlyeq$, car}
\[
\deleted{\unsure{G|F}\preccurlyeq F\Longleftrightarrow G\ll F}
\]

b) Chacune des relations $\leq$, $\ll$ détermine la relation $\parallel$. En effet
\[
\left\{\begin{aligned}
\text{a)}\quad F\parallel G &\ \deleted{\unsure{\Longleftrightarrow|\Longrightarrow}}\unsure{\Longrightarrow|\Longleftrightarrow}\
 \{F,G\}\ \text{majoré pour }\leq,\ \text{et }\mathop{\mathrm{Inf}}\limits_{\leq}(F,G)=\varnothing_F\\
\text{b)}\quad F\parallel G &\ \Longleftrightarrow\
 \deleted{\unsure{\mathrm{Inf}|\mathrm{S}}}\mathop{\mathrm{Sup}}\limits_{\ll}(F,G)\ \text{existe, et }\mathop{\mathrm{Inf}}\limits_{\ll} F,G=\varnothing_F
\end{aligned}\right.
\tag*{(1.24)}
\]
\margin{\unsure{Verifier|Vérifier} $\longrightarrow$}
(\uline{NB} $\varnothing_{\unsure{\mathcal{F}|\mathscr{F}}}$ est à la fois plus petit \unsure{élémt|élément} de $\mathcal{F}$ pour $\leq$, et pour $\ll$)

c) La relation $\ll$ détermine $\preccurlyeq$, par
\margin{$\times$}
\[
\begin{aligned}
F\preccurlyeq G &\Longleftrightarrow F\ll G,\ \text{et }\forall H\underset{\deleted{\unsure{\leq|<}}}{\in}\mathcal{F}\deleted{\unsure{K|H}}\ \text{on a}\\
&\hspace{2em}(H\parallel F\Longleftrightarrow H\parallel G)
\end{aligned}
\tag*{(1.25)}
\]
et on applique b).

d) La relation $\ll$ détermine $\mathcal{L}\subset\mathcal{F}$ comme l'\unsure{ens} des éléments \unsure{minimaux} de $\mathcal{F}\setminus\{\varnothing_{\mathcal{F}}\}$ :
\[
\mathcal{L}=\bigl\{\substack{X\in\mathcal{F}\\ X\neq\varnothing_{\mathcal{F}}}\mid
 Y\in\mathcal{F},\ Y\ll X,\ Y\neq\varnothing_{\unsure{\mathcal{F}|F|f}}\ \Rightarrow Y=X\bigr\} \tag*{(1.26)}
\]
(un \unsure{lieu|lien} est une figure non vide dont le seul raffinement non vide est elle-$\overline{\mathrm m}$).

\deleted{Ceci dit, $\ll$ implique la relation $\mathrel{\tikz[baseline=0pt,x=1em,y=1em]{\draw (-.38,.62)--(.38,.25)--(-.43,-.10); \draw (.34,.62)--(-.32,.25)--(.38,-.10);}}$ par}
\[
\tikz[baseline=(trial.base)]{
\node[inner xsep=5pt,inner ysep=6pt] (trial) {$\deleted{\unsure{F|G}\mathrel{\tikz[baseline=0pt,x=1em,y=1em]{\draw (-.38,.62)--(.38,.25)--(-.43,-.10); \draw (.34,.62)--(-.32,.25)--(.38,-.10);}} G\Longleftrightarrow\{\unsure{\forall|V}}$};
\draw (trial.north west) rectangle (trial.south east);
\draw ([xshift=.2em]trial.south west) -- ([xshift=1.45em]trial.north west);
\draw ([xshift=2.35em]trial.south west) -- ([xshift=3.6em]trial.north west);
\draw ([xshift=-1.55em]trial.south east) -- ([xshift=-.3em]trial.north east);
} \tag*{\deleted{(1.27)}}
\]
Je ne vois pas que $\ll$ ($+\ \preccurlyeq$, c'est kif-kif) détermine $\leq$, ni même $\mathrel{\tikz[baseline=0pt,x=1em,y=1em]{\draw (-.38,.62)--(.38,.25)--(-.43,-.10); \draw (.34,.62)--(-.32,.25)--(.38,-.10);}}$. Mais

\textcircled{\unsure{e|$\ominus$}}\quad $\ll$ et $\mathrel{\tikz[baseline=0pt,x=1em,y=1em]{\draw (-.38,.62)--(.38,.25)--(-.43,-.10); \draw (.34,.62)--(-.32,.25)--(.38,-.10);}}$ déterminent $\leq$ par
\margin{$\times$}
\[
G\leq F\quad\unsure{\Longleftarrow|\Longleftrightarrow}\quad (G\ll F\ \text{et}\ G\mathrel{\tikz[baseline=0pt,x=1em,y=1em]{\draw (-.38,.62)--(.38,.25)--(-.43,-.10); \draw (.34,.62)--(-.32,.25)--(.38,-.10);}} F) \tag*{(1.27)}
\]

Donc finalement,
\scanpage{14}\mspage{13}
on trouve comme possibilités \unsure{optima|optimes} de \unsure{definir|définir} la structure d'atelier sur $\unsure{\mathcal{F}|\mathscr{F}}$ :
\[
\left\{\begin{aligned}
 &\text{Soit par}\quad\boxed{\leq,\quad\ll}\\
 &\text{soit par}\quad\leq,\quad\preccurlyeq\\
 &\text{soit par}\quad\ll,\quad\mathrel{\tikz[baseline=0pt,x=1em,y=1em]{\draw (-.38,.62)--(.38,.25)--(-.43,-.10); \draw (.34,.62)--(-.32,.25)--(.38,-.10);}}\ .
\end{aligned}\right.
\tag*{(1.28)}
\]
\margin{Non, \unsure{sa|ça} ne \unsure{suffira|suffit} pas, d'après nouveau point de vue \ldots}

Dans la première mouture, j'ai choisi $\leq$, $\preccurlyeq$, en définissant en termes de ceux-ci
\[
\left\{\begin{aligned}
 &\ll\ \text{est la relation de préordre engendrée par}\\
 &\qquad\leq,\preccurlyeq,\quad \text{ou encore}\\
 &G\ll F\quad\Longleftrightarrow\quad\exists F'\ \text{avec}\quad G\leq F'\preccurlyeq F
\end{aligned}\right.
\tag*{(1.29)}
\]
% slanted
\margin{Je \unsure{changerai|changeai} le \unsure{pt} de vue page 20, et ne \unsure{supposerai} plus que $\ll$ soit déterminé ainsi en termes de $\leq,\preccurlyeq$, mais plutôt que $\preccurlyeq$ est déterminé par $\ll,\leq$}

Cela correspond au centre de gravité \deleted{\unsure{conceptuel}} \inserted{géométrique} et intuitif, qui est bien sur la notion de subdivision, plutôt que sur celle de raffinement. Mais techniquement c'était un peu laborieux, de déduire les propriétés de transitivité et \inserted{surtout} de réflexivité pour $\ll$, cela amenait à mettre des axiomes en surnombre. Toujours techniquement, il semblerait que ce soit de partir de $\leq,\ll$ qui soit le plus simple. C'est l'approche que je vais suivre \unsure{tantôt}. Quant à la possibilité de passer par $\ll$, $\mathrel{\tikz[baseline=0pt,x=1em,y=1em]{\draw (-.38,.62)--(.38,.25)--(-.43,-.10); \draw (.34,.62)--(-.32,.25)--(.38,-.10);}}$, c'est plutôt une curiosité qu'autre chose,
\scanpage{15}\mspage{14}
En résumé, en termes de $\unsure{\mathcal{F}|\mathscr{F}}$, c'est de la donnée de deux \uline{relations} \uline{d'ordre}
\[
\boxed{\underset{\overbrace{\scriptstyle\text{rel. d'ordre dans }\mathcal{F}}}{\leq,\quad\ll}\quad\text{telles que}\quad F\leq G \Rightarrow F\ll G}
\tag*{(1.30)}
\]
que je vais partir.

En termes de celles-ci, on \unsure{definit|définit} \deleted{$\unsure{M|\mathcal{F}}$}
\[
M\subset\mathcal{F}
\]
par (1.20), avec les relations $\underset{M}{\leq},\ \underset{M}{\ll}$ induites par $\underset{\mathcal{F}}{\leq},\ \underset{\unsure{\mathcal{F}|\overline{f}}}{\ll}$, puis
\[
\mathcal{L}\subset M
\]
par (1.2\unsure{6|5}). L'application
\[
\mathcal{F}\hookrightarrow\unsure{\mathcal{P}|\mathrm{P}}(M)\qquad
\deleted{\unsure{X}}F\longmapsto\widetilde F=\unsure{\mathrm{Dep}|\mathrm{Def}}(F)
\tag*{($*$)}
\]
est \inserted{alors} \unsure{définie|definie} par
\[
\widetilde F=\{X\in M\mid X\leq F\}
\tag*{(1.31)}
\]
et
\[
\mathcal{F}\hookrightarrow\mathrm{Fig}(M)\qquad
\deleted{\unsure{X}}F\longmapsto\mathrm{Fig}_{M}(F)
\tag*{($**$)}
\]
par (1.4) :
\[
\mathrm{Fig}_{M}(F)=\{\widetilde X\mid X\in\widetilde F\}\subset\mathcal{P}(M)
\tag*{\deleted{(1.32)}}
\]
avec
\[
\widetilde X=\{Y\in M\mid Y\leq X\}\quad.
\]

Ainsi, $\mathcal{F}$ est interprété comme un ens. de parties dans $M$ via ($*$), et $\leq$ \unsure{s'interprète|s'interprête} simplement par l'inclusion
\[
F\leq G\Longleftrightarrow\widetilde F\subset\widetilde G
\tag*{(1.32)}
\]
\scanpage{16}\mspage{15}
par contre la relation $\ll$ n'a pas d'interprétation simple directe, ni en termes des \inserted{\unsure{déployements|déploiements}} $\widetilde F\subset M$, ni en termes des \unsure{multidéployement|multidéployements|multidéploiement|multidéploiements} $\mathrm{Fig}(F)$; ainsi $G\ll F$ n'implique nullement $\mathrm{Fig}_{M}(G)\ll\mathrm{Fig}_{M}(F)$. Aussi l'application ($**$) et la notation $=\mathrm{Fig}_{M}(F)$ sont-ils finalement un peu bidon !

En termes de $M$, la structure est donnée par
\[
\boxed{\begin{aligned}
&\unsure{\mathcal{F}|\mathscr{F}}\subset\unsure{\mathcal{P}|\mathrm{P}}(M)\quad\text{ens de parties de }M\\
&\ll\quad\text{relation d'ordre dans }\mathcal{F},\ \text{\uline{impliquée}}\\
&\qquad\text{par l'inclusion }(\deleted{\unsure{F}}\text{notée }F\leq G)
\end{aligned}}
\tag*{(1.33)}
\]
Ça a l'air moins sympathique du point de vue simplicité des données et d'autant plus qu'il n'y a pas \inserted{à première vue} d'interprétation directe de $\ll$ en termes de \deleted{\unsure{notions}\ \unsure{de structure|de théorie}} \inserted{notions ensemblistes} familières. Mais c'est ici que vont servir les (grandes) \unsure{ombres} !

En termes de la \inserted{première} donnée (1.33\unsure{.}\unsure{)|]} \margin{(et avant d'introduire $\ll$)}, le \deleted{\unsure{tout}} premier axiome est \deleted{\illegible{1}} le fait que $\mathcal{F}$ est \deleted{une figure} \inserted{une famille d'}ens. \inserted{\uuline{pré}admissibles} \deleted{\unsure{dans}} $M$ de
% slanted
\margin{attention au ``\uuline{pré}'' !}
\deleted{\unsure{Axiome}} support $M$, \deleted{\unsure{d'où}} donc la donnée de
% slanted
\margin{\deleted{\unsure{un ens.|un ensemble} de parties \unsure{fermées} de $\unsure{F|\mathcal{F}|\lbrace}$}}
$\mathcal{F}$ \deleted{\unsure{revient à|revient} la donnée d'} \inserted{implique} une relation de préordre dans $M$ (ce sera $\leq$)\unsure{$\big/\!6\!\big/$} -- lequel sera
\scanpage{17}\mspage{16}
une relation d'ordre ; en d'autres termes

(i) $\forall x\in M$, dans l'ens ${\unsure{\mathcal{F}|\mathscr{F}}}^{x}$ des $F\in\mathcal{F}$ telles que $\unsure{X|x}\in F$, il y a un plus petit \unsure{élément}, soit $F_{x}$ \unsure{[|(}NB ce sera l'ens des $Y\in M$ tels que $Y\leq x$)

(ii) L'application $X\unsure{\mathrel{\text{---}}|\longrightarrow} F_{X}$ est injective, i.e. si $X,Y\in M$, $X\ne Y$, alors $\exists F\in\mathcal{F}$ telle que $X\in F,\ Y\notin F$, ou l'inverse $Y\in F,\ X\notin F$.

Dès lors, les $F\in\mathcal{F}$ \deleted{\unsure{sont des}} \inserted{forment un ens} de parties fermées de $M$, \inserted{contenant les $M_{\leq x}$, $x\in M$\unsure{!?|!|?}} \deleted{\unsure{Se donner un} $\mathcal{F}\subset\unsure{\mathcal{P}|\mathrm{P}}(M)$ \unsure{satisfais.|satisfaisant} (i) (ii), \unsure{revient au même que de se donner}} \inserted{De plus, on suppose} $\varnothing$ (partie vide de $M$) $\in\mathcal{F}$ (ce qui implique que $\mathcal{F}$ n'est \uline{jamais} admissible !).

Dans GF $\overline{\underline{\mathrm{IV}}}$ (Analysis Situs, \unsure{1ère|prem.} mouture) j'ai explicité en long et en large la relation entre la donnée de $(\mathcal{F},\leq)$, \unsure{satisfais.|satisfais} à des conditions qui y étaient \unsure{appelées|appellées} $\underline{\mathrm{C1}}$, $\underline{\mathrm{C2}}$ $\underline{\mathrm{C3}}$ (\unsure{pages|page} \deleted{\unsure{52|53}}\unsure{53|52} \deleted{\unsure{56|54}}\unsure{54|56}, 58),
% slanted
\margin{J'appellerai ces conditions \unsure{plutôt} $\underline{\mathrm{At}}\,1$, $\underline{\mathrm{At}}\,2$, $\underline{\mathrm{At}}\,3$}
\deleted{\unsure{que}} et un ens ordonné $M$, muni d'un ens de parties fermées $\mathcal{F}$, contenant les $M_{\leq \deleted{\unsure{x|y}}x}$, la partie vide, \inserted{et} \unsure{avec} \unsure{tout|toute} $F$ ses parties fermées (pour l'ordre de $M$), plus la condition : si $F,G\in\unsure{\mathcal{J}|\mathcal{F}}$, tels que $\forall x\in F,\ \unsure{y|\forall y}\in G$ on ait $\unsure{\overset{\scriptscriptstyle\cdots}{\scriptscriptstyle v}|\overset{\scriptscriptstyle\cdots}{\scriptscriptstyle\vee}}\,M_{\leq x}\cup M_{\leq y}\in\mathcal{F}$, alors $F\cup G\in\mathcal{F}$. (critère de compatibilité).
